\begin{tikzpicture}[
  font=\scriptsize,
  lapis/.style={draw, rounded corners, align=left, text width=7.9cm, minimum height=0.88cm},
  induk/.style={draw, rounded corners, align=center, text width=2.6cm, minimum height=2.2cm},
  panah/.style={-{Stealth[length=2mm]}, thick}]
  \node[lapis, fill=orange!12] (a) at (0,0)
    {\textbf{Mengikat.} UU 27/2022 Pasal 10 dan Pasal 34; PP 28/2024 Pasal 966 dan Pasal 972;
    Permenkes 24/2022};
  \node[lapis, fill=blue!8] (b) at (0,-1.08)
    {\textbf{Pedoman.} Surat Edaran 9/2023 tentang etika; panduan WHO untuk AI kesehatan};
  \node[lapis, fill=gray!35] (c) at (0,-2.16)
    {\textbf{Belum berlaku, dipantau.} Rancangan Perpres etika dan peta jalan AI};
  \node[lapis, fill=blue!8] (d) at (0,-3.24)
    {\textbf{Pembanding.} EU AI Act, menjadi kewajiban hanya bila cakupan Article 2 terpenuhi};
  \node[lapis, fill=blue!8] (e) at (0,-4.32)
    {\textbf{Sukarela.} ISO/IEC 42001, ISO/IEC 42005, dan NIST AI RMF};
  \node[induk, fill=green!10] (k) at (6.1,-2.16)
    {\textbf{Kebijakan KB-09}\\Kewajiban hukum, pilihan internal, dan butir pemantauan ditulis
    terpisah};
  \draw[panah] (a.east) -- (k);
  \draw[panah] (b.east) -- (k);
  \draw[panah, dashed] (c.east) -- (k);
  \draw[panah] (d.east) -- (k);
  \draw[panah] (e.east) -- (k);
\end{tikzpicture}
